\documentclass[10pt]{ctexart}
\usepackage[a4paper,margin=1in]{geometry}
\usepackage{amsmath,amssymb,booktabs}
\usepackage{enumitem}
\usepackage{microtype}
\usepackage{hyperref}
\setlist[itemize]{leftmargin=1.5em}
\title{007：Baseline 与 Step 信息使用说明}
\author{}
\date{}
\begin{document}
\maketitle
\vspace{-1.2em}

\section*{核心判断}
因为 StepPPO 显式使用 step reward、step ordering 和 step value function，所以最相关的 baseline 是 GiGPO，而不只是 GRPO。GiGPO 也是 step-aware，但没有 learned critic。因此关键问题是：shared-backbone learned critic 是否能优于 GiGPO 的 grouped relative step signal。

\section*{方法关系}
\begin{center}
\begin{tabular}{llll}
\toprule
方法 & 使用 step reward/order & learned critic & 主要角色 \\
\midrule
GRPO & 否 & 否 & non-step baseline \\
GiGPO & 是 & 否 & step-aware baseline \\
StepPPO & 是 & 是 & learned step-aware method \\
\bottomrule
\end{tabular}
\end{center}
GRPO 使用 episode-level grouped relative advantage。GiGPO 增加 step-level grouped relative advantage。StepPPO 保留 episode-level relative signal，但用 learned value-based GAE correction 表示 step signal。

\section*{什么叫不使用 Step 信息}
如果一个方法忽略 immediate step reward、step order 和 anchor observation，只使用 episode-level outcome，那么它就是 non-step 方法。即使工程上把 episode flatten 成 step samples，只要每个 step 都拿同一个 episode-level signal，它仍然属于 non-step baseline。按这个定义，GRPO 是最干净的 non-step baseline。

\section*{为什么 Classic PPO 不是干净 Baseline}
classic PPO 加 critic 后可能在 token/sample 维度上学习 value，这在 agentic setting 中会不严格地引入部分 step 结构，同时又不满足我们定义的 environment-step value。因此它不如 GRPO 适合作为 no-step 对照。

\section*{推荐实验矩阵}
\begin{itemize}
\item Non-step baseline：GRPO。
\item Step-aware critic-free baseline：GiGPO。
\item Learned step-aware method：StepPPO-v1 和 StepPPO-v0。
\item Ablation：no step term、detached value backbone、lower value coefficient、mixing 后 final-advantage normalization。
\end{itemize}

\section*{No-Step-Advantage Ablation}
直接 ablation 是设置 $w_{\mathrm{step}}=0$：
\[
A_t^{\mathrm{final}}=A_t^{\mathrm{epi}}.
\]
如果仍训练 value head 但 policy 不使用 step term，可以隔离 value regression 对 shared backbone 的干扰；如果完全关闭 value head，则更接近 GRPO 数据路径。

\section*{结论}
GRPO 用于衡量 step information 是否有收益；GiGPO 用于判断 StepPPO 的 learned critic 是否优于已有 step-aware critic-free 方法。不能只用 GRPO 作为 StepPPO 的主要 baseline。

\end{document}
